
\subsubsection{5.1 Primary Robustness Comparison (h-e1)}

Table 1 reports Spearman rho at each permutation severity level from the primary two-encoder comparison (50 training epochs, seed 42).

\textbf{Table 1: Primary robustness comparison (h-e1)}

\begin{table*}[htbp]
\centering
\resizebox{\dimexpr 2\columnwidth+\columnsep\relax}{!}{%
\begin{tabular}{lllllll}
\toprule
Encoder & rho(s=0) & rho(s=0.25) & rho(s=0.5) & rho(s=1.0) & delta-rho & Bootstrap p \\
\midrule
NFT-base & 0.4886 & 0.4886 & 0.4886 & 0.4886 & 4.09 x $10^-6$ & 0.477 \\
flat-MLP & 0.3029 & 0.2704 & 0.1945 & 0.1434 & 0.1595 & 0.000 \\
\bottomrule
\end{tabular}
}
\end{table*}

NFT-base showed no statistically significant degradation under permutation stress (p = 0.477; the null hypothesis of delta-rho = 0 was not rejected). Flat-MLP showed statistically significant degradation (p = 0.000), with a 52.7\% relative drop in Spearman rho from s=0 to s=1.0.

NFT-base achieved higher baseline Spearman rho (0.489 vs. 0.303 at s=0) despite having 40x fewer parameters than flat-MLP (75K vs. 3.04M). This baseline performance difference was not the primary hypothesis and is discussed further in Section 6.

\begin{figure}[htbp]
\centering
\includegraphics[width=\columnwidth]{figures/fig_e1_2_rho_vs_severity.png}
\caption{Figure 1: Spearman rho vs. permutation severity}
\label{fig:fig_e1_2_rho_vs_severity}
\end{figure}

\textit{Figure 1: Spearman rho as a function of permutation severity for NFT-base and flat-MLP (h-e1). NFT-base maintains rho = 0.4886 across all severity levels. Flat-MLP declines from 0.303 to 0.143.}

\begin{figure}[htbp]
\centering
\includegraphics[width=\columnwidth]{figures/fig_e1_1_delta_rho_bar.png}
\caption{Figure 2: Delta-rho bar chart}
\label{fig:fig_e1_1_delta_rho_bar}
\end{figure}

\textit{Figure 2: Permutation sensitivity (delta-rho) for NFT-base and flat-MLP. Dashed line indicates the pre-specified 0.02 threshold.}

\subsubsection{5.2 Six-Encoder Ablation and Mediation Analysis (h-m1)}

Table 2 reports results from the six-encoder ablation (100 training epochs, 3 seeds, 18 total runs).

\textbf{Table 2: Six-encoder ablation results (h-m1, mean across 3 seeds)}

\begin{table*}[htbp]
\centering
\resizebox{\dimexpr 2\columnwidth+\columnsep\relax}{!}{%
\begin{tabular}{lllllll}
\toprule
Encoder & rho(s=0) & rho(s=1.0) & delta-rho & R-squared(s=0) & p (Holm) & Significant \\
\midrule
Oracle-canon & 0.465 & 0.465 & 0.000 & 0.216 & 1.0 & No \\
NFT-base & 0.489 & 0.489 & 4.71 x $10^-7$ & 0.300 & 0.829 & No \\
NFT+aug & 0.489 & 0.489 & 2.32 x $10^-7$ & 0.300 & 1.0 & No \\
flat-MLP+aug & 0.237 & 0.014 & 0.224 & 0.072 & 0.000 & Yes \\
flat-MLP & 0.303 & -0.337 & 0.640 & 0.092 & 0.000 & Yes \\
flat-MLP+canon & N/A & N/A & NaN & N/A & 0.000 & N/A \\
\bottomrule
\end{tabular}
}
\end{table*}

Flat-MLP+canon collapsed to constant predictions (output std approximately 0 across all 3 seeds) and is reported as N/A.

The flat-MLP Spearman rho reversed sign under full permutation (from +0.303 to -0.337), indicating that the encoder's learned representation became anti-correlated with the target under permuted inputs. This effect was larger than in h-e1 (delta-rho = 0.640 vs. 0.160), consistent with the longer training (100 vs. 50 epochs) potentially increasing reliance on neuron ordering statistics.

\textbf{Mediation analysis:} delta-R-squared = R-squared(NFT-base) - R-squared(flat-MLP+aug) = 0.300 - 0.072 = 0.228. This exceeded the pre-specified 0.10 threshold by a factor of 2.28. The result indicates that NFT embeddings explained 22.8 additional percentage points of variance in generalization gap prediction compared to augmentation-based embeddings.

\begin{figure}[htbp]
\centering
\includegraphics[width=\columnwidth]{figures/fig_m1_3_mediation_bar.png}
\caption{Figure 3: R-squared bar chart showing the mediation gap}
\label{fig:fig_m1_3_mediation_bar}
\end{figure}

\textit{Figure 3: R-squared values for each encoder at s=0. The delta-R-squared = 0.228 gap between NFT-base and flat-MLP+aug is annotated.}

\begin{figure}[htbp]
\centering
\includegraphics[width=\columnwidth]{figures/fig_m1_4_rho_heatmap.png}
\caption{Figure 4: Rho heatmap across 6 encoders and 4 severity levels}
\label{fig:fig_m1_4_rho_heatmap}
\end{figure}

\textit{Figure 4: Spearman rho for six encoders at four permutation severity levels. NFT-base and NFT+aug rows are uniform. Flat-MLP degrades with increasing severity. Flat-MLP+canon is absent (constant predictor).}

\subsubsection{5.3 Alternative Approaches (h-m2)}

The h-m2 experiment evaluated augmentation and canonicalization as alternatives to architectural equivariance, reusing h-m1 checkpoints. The pre-specified SHOULD\_WORK gate required: (a) augmentation delta-rho > 0.05, (b) canonicalization delta-rho > 0.03, (c) NFT delta-rho < 0.02, and (d) a strict ranking NFT < canon < aug < flat-MLP. The gate failed due to canonicalization collapse and the resulting inability to evaluate the ranking condition.

\textbf{Augmentation.} Flat-MLP+aug achieved mean delta-rho = 0.207, representing an approximately 67\% reduction compared to flat-MLP (delta-rho = 0.627). However, per-seed results showed high variance: delta-rho of 0.096 (seed 42), 0.210 (seed 123), and 0.317 (seed 456). The coefficient of variation was approximately 107\%, indicating that the augmentation benefit was strongly seed-dependent. The bootstrap test confirmed that flat-MLP+aug remained significantly less robust than NFT-base (p = 0.000, Holm-corrected).

\textbf{Table 3: Per-seed augmentation results (h-m2)}

\begin{table*}[htbp]
\centering
\resizebox{\dimexpr 2\columnwidth+\columnsep\relax}{!}{%
\begin{tabular}{lllll}
\toprule
Encoder & Seed 42 delta-rho & Seed 123 delta-rho & Seed 456 delta-rho & Mean delta-rho \\
\midrule
flat-MLP & 0.649 & 0.605 & 0.626 & 0.627 \\
flat-MLP+aug & 0.096 & 0.210 & 0.317 & 0.207 \\
flat-MLP+canon & NaN & NaN & NaN & NaN \\
NFT-base & 2.86 x $10^-7$ & 2.81 x $10^-7$ & 1.75 x $10^-7$ & 2.47 x $10^-7$ \\
\bottomrule
\end{tabular}
}
\end{table*}

\textbf{L2-norm canonicalization.} Flat-MLP+canon produced constant predictions across all 3 seeds (output std approximately 0, all predictions approximately 0.0006). Root cause analysis attributed this to L2 normalization projecting all weight vectors onto the unit sphere, thereby destroying the relative magnitude information that distinguishes networks with different generalization gaps. This represents a systematic failure of this canonicalization approach for this prediction task.

\textbf{Oracle canonicalization.} Oracle-canon achieved delta-rho = 0.000, confirming the theoretical upper bound for post-hoc canonicalization with oracle access to the reference neuron ordering.

\begin{figure}[htbp]
\centering
\includegraphics[width=\columnwidth]{figures/fig_m1_1_delta_rho_bar.png}
\caption{Figure 5: Six-encoder delta-rho bar chart}
\label{fig:fig_m1_1_delta_rho_bar}
\end{figure}

\textit{Figure 5: Permutation sensitivity (delta-rho) for all six encoders. Dashed line indicates the 0.02 threshold. NFT-base and NFT+aug fall below the threshold; flat-MLP variants fall above.}

\subsubsection{5.4 Parameter Efficiency}

\textbf{Table 4: Parameter count and baseline performance}

\begin{table}[htbp]
\centering
\resizebox{\columnwidth}{!}{%
\begin{tabular}{llll}
\toprule
Encoder & Parameters & rho(s=0) & Final Train Loss \\
\midrule
NFT-base & 75K & 0.489 & 5.0 x $10^-5$ \\
flat-MLP & 3,040K & 0.303 & 6.3 x $10^-5$ \\
flat-MLP+aug & 3,040K & 0.237 & not specified \\
\bottomrule
\end{tabular}
}
\end{table}

NFT-base achieved higher Spearman rho at baseline (0.489 vs. 0.303) with 40x fewer parameters. Both models converged to similar final training losses (5.0 x $10^-5$ vs. 6.3 x $10^-5$). This comparison is confounded by unmatched parameter counts: the performance difference could reflect structural inductive bias, the flat-MLP being over-parameterized relative to the dataset size (3.04M parameters / 23,997 training examples), or both. A matched-parameter comparison was not conducted.
